% Direct-authority diplomatic French for printed p.189.
\clearpage
\SourcePage{189}

\LayerTarget{tohoku:eq:4.2.8}
\noindent (qui sont caractérisés par la condition de définir un homomorphisme
de foncteurs cohomologiques, se réduisant à l'homomorphisme (4.1.3) en
dimension 0).
\LayerTarget{tohoku:proof:theorem-4.2.2}
$A$ étant fixé, pour que ces homomorphismes soient des isomorphismes quel que
soit $p$ et $B$, il faut et il suffit que (i) ce soit le cas pour $p=0$
(ii) $\mathrm{Ext}_{{\mathbf O}(x)}^p(A(x),B(x))=0$ si $p>0$ et $B$ est un ${\mathbf O}$-module
injectif. Les propositions 4.1.1 et 4.1.2 ont visiblement été faites exprès
pour vérifier ces conditions. On obtient par suite :

\LayerTarget{tohoku:theorem:4.2.2}
\textsc{Théorème 4.2.2.} \textit{Supposons que ${\mathbf O}$ soit un faisceau
cohérent d'anneaux noethériens à gauche, et que $A$ soit un ${\mathbf O}$-module
cohérent. Alors les homomorphismes \textup{(4.2.8)} sont des isomorphismes quel que soit
$B$ et $p$.}

\LayerTarget{tohoku:remark:locally-free-ext-vanishing}
Notons enfin le cas trivial suivant où les homomorphismes (4.2.8) sont des
isomorphismes, les deux membres étant nuls :

\LayerTarget{tohoku:proposition:4.2.3}
\textsc{Proposition 4.2.3.} \textit{Supposons que $A$ soit localement
isomorphe à ${\mathbf O}^n$ \textup{(}où $n$ est un entier fini $>0$ donné\textup{)}. Alors
$\mathbf{Ext}_{\mathbf O}^p(A,B)=0$ pour tout ${\mathbf O}$-module $B$ et $p>0$.}

\LayerTarget{tohoku:proof:proposition-4.2.3}
En effet, en vertu de la nature locale de $\mathbf{Ext}_{\mathbf O}^p(A,B)$, on peut
supposer $A={\mathbf O}^n$. Mais alors $\mathbf{Hom}_{\mathbf O}(A,B)=B^n$ est un foncteur
\textit{exact} en $B$, d'où la conclusion.

\LayerTarget{tohoku:corollary:4.2.3}
\textsc{Corollaire.} \textit{Si $A$ est localement isomorphe à
${\mathbf O}^n$, on a}
\[
\mathrm{Ext}_{\mathbf O}^p(X;A,B)=H^p(X,\mathbf{Hom}_{\mathbf O}(A,B))
\]
\textit{pour tout ${\mathbf O}$-module $B$.}

\LayerTarget{tohoku:proof:corollary-4.2.3}
Cela résulte en effet aussitôt de la suite spectrale du théorème 4.2.1. En
particulier, on trouve (faisant $A={\mathbf O}$) :
\LayerTarget{tohoku:eq:4.2.9}
\[
\tag{4.2.9}
H^p(X,B)=\mathrm{Ext}_{\mathbf O}^p(X;{\mathbf O},B)
\]
\LayerTarget{tohoku:remark:cohomology-as-ext}
formule valable sans aucune hypothèse sur ${\mathbf O}$ ou le ${\mathbf O}$-module $B$, et qui est
d'ailleurs triviale directement, les $H^p(X,B)$ étant les foncteurs dérivés,
\textit{dans $\mathbf C^{\mathbf O}$} (grâce à prop. 4.1.3) du foncteur
$\Gamma(B)=\mathrm{Hom}_{\mathbf O}({\mathbf O},B)$.

\LayerTarget{tohoku:construction:computing-spectral-sequence}
Pour terminer ce N$^\circ$, nous indiquons une ``méthode de calcul'' de la
suite spectrale du théorème 4.2.1, plus commode que celle qui résulterait de
l'application pure et simple de la définition.

\LayerTarget{tohoku:proposition:4.2.4}
\textsc{Proposition 4.2.4.} \textit{Soit
$\mathbf{L}=(L_i)$ une résolution gauche de $A$ par des ${\mathbf O}$-modules $L_i$
localement isomorphes à des ${\mathbf O}^{n_i}$, considérons le foncteur
$\mathbf{h}_{\mathbf{L}}$ sur $\mathbf C^{\mathbf O}$, à valeurs dans les complexes à degrés
positifs dans $\mathbf C$, défini par
$\mathbf{h}_{\mathbf{L}}(B)=\mathbf{Hom}_{\mathbf O}(\mathbf{L},B)$; posons de même
$\mathbf{h}_A(B)=\mathbf{Hom}_{\mathbf O}(A,B)$. Alors $\mathbf{h}_{\mathbf{L}}$ est une
résolvante \textup{(}cf.\ \textup{2.5)} du foncteur $\mathbf{h}_A$.}

\LayerTarget{tohoku:proof:proposition-4.2.4}
En effet, $\mathbf{h}_{\mathbf{L}}$ est un foncteur exact (prop. 4.2.3), de
plus $H^0(\mathbf{h}_{\mathbf{L}}(B))=\mathbf{h}_A(B)$ grâce au fait que le
foncteur $C\to\mathbf{Hom}_{\mathbf O}(C,B)$ est exact à gauche en $C$, enfin si
$B$ est injectif, $\mathbf{h}_{\mathbf{L}}(B)$ est un complexe acyclique en
dimensions $>0$, puisque alors le foncteur
$C\to\mathbf{Hom}_{\mathbf O}(C,B)$ est exact.

\LayerTarget{tohoku:corollary:4.2.4:1}
\textsc{Corollaire 1.} \textit{Sous les conditions de la proposition
précédente, on a}

\TohokuContinuation{190}
